\documentclass[10pt,a4paper]{article}

	\usepackage{amsmath,amssymb,amsthm,mathtools,amsfonts}
	\usepackage{breqn}
	\usepackage[utf8]{inputenc}
	\usepackage[english]{babel}
	\usepackage[T1]{fontenc}
	\usepackage{graphicx}
	\graphicspath{{../img/}}

	%%%%%%%%%%%%%%%%%%%%%%%%%%%%%%%%%%%%%%%%%%%%%%%%%%%%%%%%%%%%%%%%%%%%%%%%%%%%%% Fonts

	\usepackage{fontspec}
	\setmainfont{Times New Roman}
	\newfontfamily\garamond[Numbers=OldStyle]{EB Garamond}
	\newfontfamily\plantin[Numbers=OldStyle]{Plantin MT Pro}
	\newfontfamily\wal{Walbaum Com}
	\newfontfamily\quotefont[Ligatures=TeX]{Linux Libertine O}
	\newfontfamily\vest{TRY Vesterbro Regular}
	\newfontfamily\vestb{TRY Vesterbro Poster}
	\newfontfamily\overp{Overpass}

	\newfontfamily\mysectionfont[Numbers=OldStyle]{Klavika light}

	\usepackage[pdfborder={0 0 0}]{hyperref} %% Ingen firkant rundt om referencer
	\usepackage{multicol}
	\usepackage{tabularx}
	\usepackage{enumitem}
	\usepackage{wrapfig}
		\setlength{\columnsep}{0pt}
		\setlength{\intextsep}{-10pt}
	\usepackage{caption}
	\usepackage[svgnames]{xcolor}
	\usepackage{dirtytalk}
	\usepackage[modulo]{lineno}
	\usepackage{hyperref}
	\usepackage{pdfpages}
	\usepackage{tasks}
	\newcommand\svar[1]{\newline\textcolor{red}{\textbf{Svar}\\#1}}
	\newcommand{\qm}[1]{"#1"}

	%%%%%%%%%%%%%%%%%%%%%%%%%%%%%%%%%%%%%%%%%%%%%%%%%%%%%%%%%%%%%%%%%%%%%%%%%%%%%% Colors

	\usepackage{xcolor}
  \usepackage{transparent}

	\definecolor{frbblue}{RGB}{47, 88, 109}
	\definecolor{hgred}{RGB}{140, 10, 10}
	\definecolor{frbgreenl}{RGB}{162, 202, 137}
	\definecolor{frbgreen}{RGB}{28, 85, 66}
	\definecolor{frbgrey}{RGB}{243, 245, 242}


	\usepackage{titlesec}
	\titleformat{\subsection}{\color{frbgreen}\overp\large}{}{}{}
	\newcommand{\source}[1]{Source: \href{#1}{#1}}
	\usepackage{lipsum}
	\usepackage{comment}
	\usepackage{ulem}

	%\reversemarginpar
	\newcommand{\glos}[3]{\dotuline{#1}\marginpar{\scriptsize\textit{#3} #2}}

	\setlength{\parskip}{0.5\baselineskip}
	\setlength{\parindent}{0cm}
	\setlist[enumerate]{label=\alph*),left=20pt}

	%Titling
	\usepackage{titling}
	\setlength{\droptitle}{-12ex}
	\pretitle{\begin{center}\vestb\Huge\color{frbgreen}}
	\posttitle{\end{center}}
	\preauthor{\begin{center}\vest\Large\color{frbblue}}
	\postauthor{\end{center}}
	\predate{\begin{center}\vest}
	\postdate{\end{center}}

	\setcounter{secnumdepth}{0}

	\usepackage{tikz}

\begin{document}

%\pagecolor{FloralWhite}
\title{Eksamen 6l26en}
%\date{\vspace{-5em}} % I tilfældet ingen dato
\date{Sommer, 2026}
\author{Spørgsmål 5}
\maketitle

\begin{tikzpicture}[remember picture,overlay]
    \node[anchor=north east,yshift=-15pt,xshift=-45pt]%
       at (current page.north east)
        {
\includegraphics[scale=0.1]{frbvuc_logo.png}};
    \end{tikzpicture}

\thispagestyle{empty}

\vspace{-4em}

\begin{center}
{\color{frbblue}\rule{1\textwidth}{1pt}}
\end{center}

\subsection{Prøvemateriale}
\begin{description}
\item[Theme:] Borders
\item[Text:] Sandy Rivera, "Undocumented, Yet Hopeful" (2018) (speech)
\item[Instruction:] You must make a presentation (roughly 6 minutes) of the text and relate it to the theme. In your presentation, you must
  \begin{itemize}
\item sum up the main points of the speech
\item and make an analysis of the speaker's rhetorical strategies and intention.
\end{itemize}.
\end{description}

\subsection{Bedømmelseskriterier, hf C}
Ved den \textbf{mundtlige prøve} lægges der vægt på, at eksaminanden
\begin{itemize}
  \item behersker et forståeligt og sammenhængende engelsk
  \item præsenterer og perspektiverer prøvematerialet
  \item foretager grundlæggende analyse af prøvematerialet med anvendelse af fagets tekstanalytiske grundbegreber
  \item anvender den viden, der er opnået i arbejdet med det studerede emne.
\end{itemize}
Der lægges i bedømmelsen vægt på, at eksaminanden kan indgå i uddybende samtale om præsentationen.

Der gives én karakter ud fra en helhedsvurdering af eksaminandens præstation.


\includepdf[pages=1-]{Spørgsmål_5_tekst.pdf}


\end{document}
